\section{Introduction}
\IEEEPARstart{I}{ndustrial-agent} deployments must connect each task to its
production context, process signals, and equipment interfaces. Existing work
addresses end-to-end automation control, operator-reviewed control-sequence
generation, modular factory skills, and production scheduling
\cite{xia2025control,huang2026made,koehle2026llmdriven,waseem2026pretrained}.
Other studies target PLC code generation, context-aware control, resource
capability discovery, shop-floor coordination, digital-twin fault handling,
natural-language machine access, and robotic-line control
\cite{liu2026agents4plc,wu2025instructmpc,lim2025llmenabled,zhao2026shopfloor,gill2025llmdt,hofmann2025nlc}.
Across production lines, agent roles, signal maps, protocols, objectives, and
product/run context vary. An integration architecture must accommodate these
differences around shared task and operation services while keeping each
action tied to its active production context.

Consider two production scenes with different process signals and interfaces:
cast-film thickness control and annealing-zone temperature control. An agent
task in either scene needs production context, typed signal access, and a
route from a named operation to the configured device protocol. The framework
must also distinguish a command receipt from a subsequent process
observation. The shared task and operation services can therefore be
configured with scene-specific signal maps, protocols, and limits while
retaining a common execution path. The evaluation instantiates this path in
two simulator scenes.

This paper presents AgentWorkShop as a reference implementation of
task-to-device integration for industrial agents. It connects collaborative tasks
and registered harnesses to production context, manufacturing execution
system (MES) data, analysis and modeling services, data-acquisition (DAQ) and
write-control (DCW) nodes, and protocol-specific drivers. Its interfaces include
command-line, web, SDK, REST, Model Context Protocol (MCP), and Agent-to-Agent
(A2A) surfaces. Task and chat routes enter the
collaboration runtime, whereas a direct REST tool-invocation route can call
the shared operation dispatcher without creating a task.
Scene blueprints define signal maps, objectives, and guards; channel profiles,
tool policies, and agent bindings scope operations. Agents and callers invoke
named operations through shared runtime checks; the framework does not choose
routes autonomously and retains proposals, command confirmation, and process
observations as distinct records. Fig.~\ref{fig:arch} presents the system
architecture, and Fig.~\ref{fig:runtime} shows the cross-module task flow.
Section~IV details the node-exploration and recipe lifecycles.

This paper asks whether one registered-harness path can execute bounded tasks
across two simulator scenes with different signal and protocol mappings while
retaining an auditable relation from each task through approval and verified
route receipt to a subsequent run-tagged measurement. The question concerns
the recorded task-to-device path; it does not test an immutable context
snapshot across prompt, approval, and dispatch.

The paper makes two contributions:
\begin{enumerate}
\item It specifies and implements a task-to-device integration path that
  joins collaborative tasks and registered harnesses to turn-time production
  context, per-call operation checks, semantic I/O, configured protocol
  drivers, and service-owned records. Its design distinction is the explicit
  separation of grounding context, dispatch authority, device receipt, and
  later process observation; it introduces no new protocol or
  process-control algorithm.
\item It evaluates that path using one registered harness in two simulator
  scenes with distinct signal and protocol mappings. Six sequential tasks
  completed with synthetic approval and verified readback, and all six
  records passed task-to-approval, receipt, and active-run sample linkage
  checks. None of the samples reached its target band. Separate regression
  tests exercise denial, expiry, limit, and post-approval revalidation cases.
\end{enumerate}

Section~II reviews related industrial-agent and integration work.
Section~III presents the design rationale, architecture, and cross-module
execution flow; Section~IV explains module responsibilities and their
integrated operation. Section~V describes the experiments and results, and
Sections~VI and~VII discuss the evidence limits and conclude the paper.
